Varias limitaciones deben matizar la interpretación de los resultados. En primer lugar, el conjunto supervisado es un subconjunto procesado y no una muestra representativa de toda la población de contrataciones. La inclusión exigía disponibilidad documental, extracción exitosa del texto, generación exitosa de embeddings y uno de los dos estados finales definidos. Como el filtrado se debió principalmente a restricciones de tiempo y presupuesto de procesamiento, las conclusiones deben entenderse como aplicables al subconjunto procesado.

En segundo lugar, el subconjunto procesado presenta menos desbalance que la población real. Los datos experimentales contienen alrededor de $73\%$ de ejemplos positivos, mientras que la población general está mucho más inclinada hacia resultados exitosos, con aproximadamente $94\%$ de positivos. Esta diferencia importa porque las métricas dependientes del umbral y las interpretaciones sensibles a la prevalencia no son invariantes al balance de clases. Los puntos operativos informados no deben considerarse directamente desplegables sin recalibrarlos en el entorno objetivo.

En tercer lugar, la falta de transferibilidad del umbral es una preocupación práctica. Los umbrales optimizados en validación mejoran algunas métricas de la muestra actual, especialmente para TF--IDF más regresión logística, pero pueden cambiar bajo otras prevalencias, contextos institucionales o supuestos sobre el costo de los errores. Por tanto, el ajuste debe interpretarse como una herramienta de análisis para este conjunto y no como una prescripción universal.

En cuarto lugar, la principal preocupación residual no es una fuga directa de información sobre el resultado final. Como los documentos se producen durante la convocatoria, tal fuga parece poco probable. Los riesgos más relevantes son las correlaciones predictivas indirectas, el aprendizaje de atajos, la representatividad limitada y el desajuste de distribución. Por ejemplo, los modelos podrían aprovechar regularidades institucionales o estilísticas predictivas en este conjunto, pero menos estables entre organismos, períodos o categorías de contratación. Los hallazgos también pueden ser específicos de la construcción vial y no transferirse sin cambios a otras categorías.
